\section{Localization of every column at the endpoint 180}\label{app:global}
This section proves the global and fine inputs to
Proposition~\ref{prop:localization}. All partitions below contain every
original column. The same partition is kept during the forward estimates.

\subsection{Deficit identity and estimates before partitioning}
Let $t_{ij}=(1+x_i\cdot x_j)/2$. The polynomial
\[
 f(t)=\alpha+\tfrac34(t-\tfrac12)-k(t^2-t+\tfrac16)
\]
satisfies $f(t)-\sqrt t=g(t)$ exactly. Put
$T_1=\sum_{ij}r_ir_jt_{ij}$. Tightness gives
$\sum_i r_i^2t_{ij}=N/2$, so
\[
 N^2/2-T_1=\tfrac12\sum_{ij}(r_i-r_j)^2t_{ij}\ge0.
\]
Also
$\sum_{ij}r_ir_j(t_{ij}^2-t_{ij}+1/6)
=\frac14\|\sum r_ix_ix_i^T-(L/3)I\|_F^2$.
Summing the polynomial identity proves
\begin{equation}\label{app:global:deficit}
 \alpha N^2-\Gamma_I(Y)=u+v+w+\tfrac34(N^2/2-T_1).
\end{equation}
Every term on the right is nonnegative. In particular $L\le N$.

Set $c=13/4$. On $N\ge180$, $u+v+w\le M\le11/19$ gives
$L/N>9999/10000$, using $a>77/250$. Since
$(77/250)(19999/10000)>8/13$,
\begin{equation}\label{app:global:radial}
 D_r:=\sum(r_i-1)^2=\frac{2u}{a(N+L)}\le\frac{cu}{N},
 \quad s=\sqrt{cu/N}<\frac{103}{1000}.
\end{equation}
Thus every radius is positive. With $E=\sum_{ij}g(t_{ij})$,
\begin{equation}\label{app:global:energy}
 E\le\frac{w}{(1-s)^2}\le\frac{M}{1-cM/N}<\frac{293}{500}.
\end{equation}
The second bound retains $u+w\le M$: after setting $x=\sqrt{cu/N}$,
$m=cM/N$, its numerator difference is
$m(1-x)^2-(m-x^2)(1-m)=(x-m)^2$.

Put $R_x=\sum x_i$, $M_0=\sum x_ix_i^T$, and
$l_i=(1,\sqrt3x_i)^T/2\in\mathbb R^4$. Cauchy--Schwarz and
$\sum r_i^2x_i=0$ give
\[
 \|M_0-NI/3\|_F\le2\sqrt{v/k}+\sqrt{2cu/3},\quad
 \|M_0-NI/3\|_{\rm op}<23/10,
\]
\[
 \|R_x\|^2\le(N/3+23/10)(2+s)^2D_r.
\]
For the operator estimate, a traceless symmetric $3\times3$ matrix
has squared operator norm at most two thirds of its squared Frobenius
norm, and
$(2/3)(1000/87+13/6)M<(23/10)^2$. The exact lifted identity is
\[
 \left\|\sum l_il_i^T-NI_4/4\right\|_F^2
 =\tfrac38\|R_x\|^2+\tfrac9{16}\|M_0-NI_3/3\|_F^2
 \le T(u,v):=Au+Bv+J\sqrt{uv},
\]
where
\begin{equation}\label{app:global:moment}
 A=\frac{19739570097}{6400000000},\quad B=\frac{375}{58},
 \quad J=\frac{1123}{200},\quad T(u,v)\le\frac{81}{10}(u+v).
\end{equation}
The last quadratic comparison has positive diagonal entries and determinant
\[
 \frac{293253687011}{928000000000}>0.
\]
If $z=u+v$, the traceless
$4\times4$ spectral inequality therefore yields
\begin{equation}\label{app:global:spectral}
 \lambda_{\min}\Bigl(\sum l_il_i^T\Bigr)
 \ge\Lambda_N(z):=N/4-\sqrt{243z/40}>0.
\end{equation}
Indeed an eigenvalue $t$ of a traceless $4\times4$ real symmetric matrix
forces the other squared eigenvalues to sum to at least $t^2/3$.

\subsection{Simultaneous anchors and coverage}
Let $\mathcal L=[l_1\ \cdots\ l_N]$ and $F=\mathcal L\mathcal L^T$.
Give each four-element subset $A$ probability
$\det(\mathcal L_A)^2/\det F$. For a set $J$ of size one or two,
\begin{equation}\label{app:global:inclusion}
 \Pr(J\subset A)=\det(\mathcal L_J^TF^{-1}\mathcal L_J).
\end{equation}
To verify this identity, whiten $U=F^{-1/2}\mathcal L$, so $UU^T=I$.
Cauchy--Binet expands
$\det(U\diag(t_i)U^T)=\sum_{|A|=4}\det(U_A)^2\prod_{i\in A}t_i$.
Taking the derivatives indexed by $J$ at $t_i=1$, and alternatively
expanding $\det(I+U\diag(t_i-1)U^T)$, gives
\eqref{app:global:inclusion}.

Write $q(t)=(3t-1)/2$, so $l_i^Tl_j=q(t_{ij})$.
Loewner order in the one- and two-point determinants gives
\[
 \Pr(i\in A)\le\Lambda_N(z)^{-1},\qquad
 \Pr(i,j\in A)\le[1-q(t_{ij})^2]\Lambda_N(z)^{-2}.
\]
The scalar bounds
$q(t)^2\le24g(t)$ for $t<2/3$ and
$1-q(t)^2\le96g(t)$ for $t\ge2/3$ follow from the canceled ratios
below, their derivatives and their values at $2/3$.
Explicitly, $g(t)/q(t)^2$ is at least its value at $2/3$, which exceeds
$1/24$, on the lower interval. On the upper interval
$g(t)/(1-t)>1/32$ and
$1-q(t)^2=(3/4)(1-t)(1+3t)\le3(1-t)$. Endpoints with zero factors
are handled directly or by continuity.
If $E_A$ is the sum of selected angular row energies, $H$ counts
selected unordered pairs with $t\ge2/3$, and
$X=2\sum_{\{i,j\}\subset A,t_{ij}<2/3}q(t_{ij})^2$, then
\[
 \mathbb E E_A\le E/\Lambda_N(z),\quad
 \mathbb E X\le24E/\Lambda_N(z)^2,\quad
 \mathbb E H\le48E/\Lambda_N(z)^2.
\]
Let $K_A=441/100$, $\eta_A=8/25$. For $E>0$, the expected potential
$E_A/(K_AE/N)+X/\eta_A^2+H$ is at most
\begin{equation}\label{app:global:potential}
 \frac{N}{K_A\Lambda_N(z)}+
 \left(24/\eta_A^2+48\right)
 \frac{(M-z)/(1-\sqrt{cz/N})^2}{\Lambda_N(z)^2}<1.
\end{equation}
The strict inequality is certified on $256$ closed intervals covering
$0\le z\le M$ at $N=180$, as specified in
Section~\ref{app:certificates}. For each interval $[l,h]$, replace
$\Lambda$ by $45-\sqrt{243h/40}$ and the energy by
\[
 \min\{M/(1-cM/180),(M-l)/(1-\sqrt{ch/180})^2\}.
\]
Every denominator is positive. This is a closed interval proof, not a
sampled maximization. Thus one actual four-subset has $H=0$, $x=\sqrt X<\eta_A$
and $E_A<(K_AE/N)(1-x^2/\eta_A^2)$.

The selected lifted Gram has smallest eigenvalue at least
$1-\sqrt3x/2$. The exact identity
\[
 \frac2{1+\sqrt{1-h}}(1-y)-(1-y^2/h)
       =\frac{(y-1+\sqrt{1-h})^2}{h}\ge0
\]
gives the coverage and disjointness ratios
\[
 R_{\rm cov}=\frac2{1+\sqrt{1-3\eta_A^2/4}},\qquad
 R_{\rm dis}=\frac2{1+\sqrt{1-\eta_A^2/2}}.
\]
For $h=\sqrt t$ and $b=1/\sqrt3$,
\[
 \frac{g(t)}{q(t)^2}=\frac{4(1-h)(3/4+kh)}{9(h+b)^2},\qquad
 \frac{g(t)}{1-t}=\frac{(h-b)^2(3/4+kh)}{1+h}.
\]
The first ratio decreases, with derivative $-4\sqrt3/[9(h+b)^3]$;
the second increases for $h\ge b$. The exact endpoint checks are
\[
 R_{\rm cov}\frac{K_A(293/500)}{180}
   <\frac{g(173/200)}{q(173/200)^2},\qquad
 R_{\rm dis}\frac{K_A(293/500)}{180}<\frac1{50}.
\]
If a data column were outside every $173/200$ cap, its unit lifted vector
would give $\sum_a q(t_{iA_a})^2\ge1-\sqrt3x/2$; its incident energies
would then contradict the first check. Two $83/100$ memberships of an
actual data column would imply an anchor squared distance at most
$(400/3)E_A$, using $g(t)/(1-t)>3/50$ there, whereas the selected-pair
bound gives at least $(8/3)(1-x/\sqrt2)$. The second check excludes it.
These assertions concern the data columns; the continuous caps need
not be disjoint. Since $173/200>83/100$, they give one partition into
four nonempty groups.

Let $\bar x_a$ be the means and
$S=\sum_a\sum_{i\in a}\|x_i-\bar x_a\|^2$. A valid lower membership
ratio at $173/200$ is
\[
 m_0=\frac{667855014484621414278044848317828318861}
 {9650268809435000000000000000000000000000}.
\]
Since means minimize squared distance,
$S\le4E_A/m_0<4K_AE/(m_0N)<255E/N$. We may weaken this to
$S\le264E/N$ for the subsequent common bootstrap.
If $E=0$, all $t_{ij}$ belong to $\{1/3,1\}$.
The positive lifted frame bound supplies four spanning data columns;
their lifted vectors are orthonormal. Every other unit lifted vector
has inner products in $\{0,1\}$ with that basis, and is one of them.
This gives the same partition and $S=0$ without division by $E$.

\subsection{Counts before fine localization}
Let $\Sigma_a$ be the within-group covariance, $S_a=\tr\Sigma_a$,
$\delta_a=m_a-N/4$, and let the square matrix $A_0$ have columns
$\sqrt{m_a}(1,\sqrt3\bar x_a)^T/2$. If
$F_0=\sum l_il_i^T-NI/4$, $H_0=\diag(0,\sum\Sigma_a)$ and
$t_0=(N-3S/4)/4$, then
\[
 A_0A_0^T-t_0I=F_0-\tfrac34(H_0-SI/4),\quad
 \diag(A_0^TA_0-t_0I)=\delta-\tfrac34(S_a-S/4)_a.
\]
The two Gram matrices have identical spectra. Their Frobenius norms,
diagonal contraction, and the centered PSD estimates imply
\begin{equation}\label{app:global:count}
 \sqrt C\le\sqrt{T(u,v)}+\frac{3\sqrt3}{4}S
                     \le\sqrt{T(u,v)}+\frac{13}{10}S.
\end{equation}
Each centered norm is at most $\sqrt3S/2$; both variance terms are
necessary. With $S\le264E/N$ and the shared budget, $256$ closed
$z$-interval checks give $C<5$ for $M=11/19$ and $C<19/4$ for $M=1/2$.
This precedes, and therefore does not rely on, any fine estimate.

For residue two the $\delta_a$ are half-integers with zero sum. The
strict square bound allows only $|\delta_a|\le3/2$, with at most one
entry of that magnitude, giving the two listed profiles. For odd $N$,
after sign reversal let $4\delta_a\equiv-1\pmod4$. Their squared sum
is below $76$. Magnitudes at least $9$ are excluded; an entry $7$
requires other entries summing to $-7$, whose minimum squares are
$25+1+1$, giving total $76$. Thus only $-5,-1,3$ remain. Their zero
sum permits either no $-5$, giving $(-1,-1,-1,3)$, or one $-5$,
giving $(-5,-1,3,3)$. Two $-5$ require two entries $3$ summing only
to $6$, so cannot have zero sum. This proves the complete count list.

\subsection{Shared-budget forward refinement}
Assume $S\le c_sE/N$ has already been proved on this partition. For a
closed region $u\le U,w\le W,u+v+w\le M$, let $\bar E$ be a valid
energy cap and set
\[
 p=1/4-m_c/N,\quad p_+=1/4+m_c/N,\quad
 S_*=c_s\bar E/N,\quad q=\sqrt{1-S_*/(pN)},\quad
 T_*=2S_* /(1+q),
\]
where $m_c=\max|\delta_a|$. All square roots and divisions require
their displayed positive domains. A sharp cap $\bar E$ follows by
maximizing $\min(W,M-u)/(1-\sqrt{cu/N})^2$ over $0\le u\le U$.
On the branch $w=M-u$ the sole interior maximum is
$u=(c/N)M^2$, of value $M/(1-cM/N)$; otherwise use the boundary.

Write
$Q=\sum_{a\ne b}m_am_b(\bar x_a\cdot\bar x_b+1/3)^2$.
The off-diagonal part of the preceding square mean-Gram calculation
gives $\sqrt Q\le(4/3)\sqrt T+S$. Write
$c_\times=\max_{a(i)\ne a(j)}(x_i\cdot x_j+1/3)$. Normalizing the
means gives the following three bounds, kept under a single budget supremum:
\begin{align}\label{app:global:finecaps}
 c_\times&\le\frac{4\sqrt T}{3Np}+b_s\sqrt E+d_sE,\\
 G&\le\frac{4\sqrt T}{3Np}+d_sE,\qquad
 \frac{\sqrt Q}{pN}\le\frac{4\sqrt T}{3Np}+d_mE,\nonumber\\
 b_s&=\sqrt{\frac{4c_s}{N(1+q)}},\quad
 d_m=\frac{c_s}{N^2p},\quad d_s=d_m\left(1+\frac{\sqrt6}{1+q}\right).
 \nonumber
\end{align}
Here $G$ bounds the Frobenius Gram error of the four normalized means.
The normalization costs follow from
$\sum_i\|x_i-\bar x_a/\|\bar x_a\|\|^2\le2S/(1+q)$.

On a closed bin $z\in[z_l,z_h]$, use
$u\le\min(U,z_h)$ and $w\le\min(W,M-z_l)$ to bound the energy and
the \emph{whole} expressions \eqref{app:global:finecaps}.
The raw moment $Au+Bv+J\sqrt{uv}$ can be maximized on $u+v=z_h$
by its concave one-variable expression. Its unconstrained upper bound
is $z_h(A+B+\sqrt{(A-B)^2+J^2})/2$; if $U_0=\min(U,z_h)$ obeys
\[
 2U_0\le z_h,\quad J^2(z_h-2U_0)^2
       \ge4(B-A)^2U_0(z_h-U_0),
\]
the constrained maximum is the value at $U_0$. This exact squared
derivative test avoids an uncertain radical branch.

There is also a spherical improvement to the cross bound. If the unit
mean dot product is at most $\gamma=-1/3+G$ and the two squared chord
errors sum to $T_s<2$, concavity of $\arccos(1-t/2)$ gives an angle
error at most $4\arcsin\sqrt{T_s/8}$. Put
\[
 C_s=1-T_s+T_s^2/8,\quad
 S_s=\sqrt{2T_s(1-T_s/8)}(1-T_s/4).
\]
When $\gamma<C_s$, the pair dot product is at most
$\gamma C_s+\sqrt{1-\gamma^2}S_s$; otherwise use $1$.
Here $T_s=2c_sE/[N(1+q)]$. This retains all within-group residuals.
On each closed budget bin, enclose both the linear bound for $c_\times$
in \eqref{app:global:finecaps} and the spherical bound just obtained,
adding $1/3$ to the latter dot-product bound. If the angular conditions
fail, use the trivial spherical bound $4/3$. Take the smaller of the
two upper enclosures on each bin, with outward rational rounding,
then the maximum over all bins covering the region. This gives a
regional upper bound. It may be reduced by taking its minimum with
other independently valid upper bounds on the same region. Denote the
resulting bound by $\bar c_\times$, so $c_\times\le\bar c_\times$
throughout the region.

Set $h_W=\sqrt{1-S_*/2}$ and
$h_C=\sqrt{(2/3+\bar c_\times)/2}$. The canceled scalar kernels are
\[
 K_W=\frac{(h_W-1/\sqrt3)^2(3/4+kh_W)}{4(1+h_W)},\quad
 K_C=\frac{(1-h_C)(3/4+kh_C)}{4(h_C+1/\sqrt3)^2}.
\]
Use lower enclosures of these kernels, $h_W>1/\sqrt3$, $h_C<1$, and
\begin{equation}\label{app:global:covariance}
 \mu_c=2\left[\frac13-\frac{23/10+S_*}{N}
                   -\frac{p_+T_*}{2(1+q)}\right],\quad
 \beta=2pK_W+\mu_cK_C>0.
\end{equation}
Then
\begin{equation}\label{app:global:jointcov}
 E\ge K_CQ+\beta NS.
\end{equation}
For detail, the cross-pair squared-dot sum is exactly
\[
 Q+2\sum_a\tr\left[\left(\sum_{b\ne a}m_b\bar x_b\bar x_b^T\right)
 \Sigma_a\right]+\sum_{a\ne b}\tr(\Sigma_a\Sigma_b).
\]
The last term is nonnegative. The full covariance is at least
$(N/3-23/10-S)I$. The variance of
$x_i\cdot(\bar x_a/\|\bar x_a\|)\in[1-T_*/2,1]$ gives its removed
same-group term at most $T_*S_a/[2(1+q)]$.
Indeed if its mean is $r_a=\|\bar x_a\|$, the scalar variance is at
most $(T_*/2)(1-r_a)$ and
$m_a(1-r_a)=S_a/(1+r_a)\le S_a/(1+q)$.
These facts yield $\mu_cNS$ for the cross part. The within-group
identity $\sum_{i,j\in a}\|x_i-x_j\|^2=2m_aS_a$ supplies
$2pK_WNS$, proving \eqref{app:global:jointcov}.

Starting from $264$, the recorded closed-domain inequalities prove
\begin{equation}\label{app:global:chain}
 264\longrightarrow240\longrightarrow183\longrightarrow107
 \longrightarrow58\longrightarrow40\longrightarrow34
 \longrightarrow32\longrightarrow31.
\end{equation}
Every arrow checks $\beta>1/c_{\rm next}$, hence is a forward implication,
not an assumed fixed point. The uniform chain uses $m_c=3/2$ and $128$
closed budget bins per stage. Subsequent regional refinement starts
from $31$, uses the already established profile's $m_c$ consistently,
and uses $64$ closed bins. No new partition is selected.

\subsection{Fine polar and spinor interface}
At the final regional coefficient let
$h_m=\sqrt{(2/3+\sqrt Q/(pN))/2}$, using its certified upper cap,
and put $K_m=K_C(h_m)$. If
$\gamma\le1/(4h_C^3)-2k$ is positive, then
$g(t)-\gamma t^2/2$ is convex on $[0,h_C^2]$. Indeed
$g^{\prime\prime}(t)=1/(4t^{3/2})-2k\ge\gamma$ for $t>0$,
and continuity extends convexity to $0$. Equivalently, $g$ is
$\gamma$-strongly convex on this interval. Jensen's inequality gives
\[
 E\ge K_mQ+[2pK_W+\mu_c\gamma/8]NS.
\]
Take a convex combination with \eqref{app:global:jointcov}, obtaining
$E\ge c_QQ+c_SNS$. If the positivity guard fails, use the first
inequality alone. The factor $1/8$ is the product of the strong-convexity
factor $1/2$ and the change of variable $t=(1+x\cdot y)/2$ squared.

Use the single polar alignment of the normalized means, with $P_G$ from
\eqref{app:candidate:polar} and $G<1/12$. The mean/residual decomposition
and Young's inequality with $\tau=201/200$ give
\[
 \sum_i\|x_i-v_{a(i)}\|^2\le A_SS+A_QQ/N,
\]
\[
 A_Q=p_+P_G\tau/p^2,\qquad
 A_S=\frac2{1+q}+
 \frac{6p_+P_G(1+1/(\tau-1))S_*}{p^2N(1+q)^2}.
\]
Therefore $L_B=\max(A_Q/c_Q,A_S/c_S)$ proves a total Bloch error
at most $L_BE/N$. This uses one joint covariance inequality, rather
than independently maximizing $Q$ and $S$.

Choose a phase floor $f$ from $127/128,63/64,31/32$ only if all three
strict guards hold:
\begin{equation}\label{app:global:phaseguards}
 L_B\bar E/N<4(1-f^2),\quad T_*<4(1-f^2),\quad
 P_GG^2<4(1-f^2).
\end{equation}
They concern, respectively, actual-to-polar, actual-to-normalized-mean
and normalized-mean-to-polar spinors. With $D_f=1/[2(1+f)]$, the
exact identity \eqref{app:candidate:spinor} and the quotient triangle
inequality give
\begin{align}\label{app:global:LU-LP}
 L_U&=D_fL_B,\\
 L_P&=\left[
 \sqrt{\frac{2D_f}{(1+q)\beta}+\frac{D_fP_G}{p^2NK_C}}
 +\frac{\sqrt{D_fP_G}\sqrt6\sqrt{\bar E/N^3}}
 {p(1+q)\beta}\right]^2.\nonumber
\end{align}
For the point estimate, the actual-to-mean contribution is bounded by
$2S/(1+q)$; the mean-to-polar contribution uses
$G\le\sqrt Q/(pN)+\sqrt6S/[pN(1+q)]$.
Maximizing their joint square-root terms under
$K_CQ+\beta NS\le E$ gives exactly the first square root in
\eqref{app:global:LU-LP}; the remaining normalization term gives the
second. Thus no separate maximum of $S$ and $Q$ is used in the first term.

In one common preliminary alignment and choice of column phases,
\begin{align}\label{app:global:physical-inputs}
 N\|Y-F\|_F^2&\le A_R:=cU+L_UW\frac{1+s_U}{(1-s_U)^2},\\
 N\max_i\|y_i-f_{a(i)}\|^2&\le P_R:=cU+L_PW\frac{1+s_U}{(1-s_U)^2},
 \qquad s_U=\sqrt{cU/N}.\nonumber
\end{align}
Indeed $d_{\rm ph}(r_iu_i,f)^2=(r_i-1)^2+r_id_{\rm ph}(u_i,f)^2$.
The sharper angular inputs are
$\sum d_{\rm ph}(u_i,f)^2\le L_U\bar E/N$ and
$\max d_{\rm ph}(u_i,f)^2\le L_P\bar E/N$.
The actual closest-gauge transport, required for the local Hessian,
is proved in Section~\ref{app:local}; a convenient preliminary
alignment is not substituted for that gauge.

\subsection{Continuous range and zero budgets}
Every accepted region is closed and all its constants cover its whole
interior and boundary. For larger $N$, at fixed deficits the energy
cap decreases, $N/\Lambda_N(z)$ decreases, and $\Lambda_N(z)$ increases.
All anchor and count inequalities therefore propagate. In the fine
chain, freeze the endpoint count floors, coefficient choices, closed
bins and phase floors. The energy and scatter decrease, mean lengths
improve, and every positive $N^{-1}$ or $N^{-1/2}$ error decreases.
The spherical cap is monotone geometrically because decreasing the
allowed mean dot product or chord error shrinks the feasible angle
range. Its algebraic maximizing branch need not remain fixed.
Propagate the finite chain in order, then its regional refinement and
polar interface. If $E=0$, \eqref{app:global:jointcov} gives $S=Q=0$;
all zero-error cases are direct. There is no division by $u,v,w$ or $E$
in the fine or local interface. These facts justify freezing the saved
180 localization coefficients for the later 182 local comparison.
